%! Sine, Cosine, and Tangent — Unit 3, Lesson 3.2 — Homework
\documentclass[10pt]{article}
\usepackage{precalculus-article}
\usepackage{precalculus-boxes}
\newcommand{\TallMath}[1]{$\displaystyle #1\rule[-1.4em]{0pt}{3.2em}$}

\begin{document}

\pageheader{Unit 3, Lesson 3.2}{Homework}
\namedateperiod

\vspace{0.10in}

\begin{notesbox}{Homework: Sine, Cosine, and Tangent}

\textbf{1.} Complete the table of exact values using the unit circle.
(Hint: use 45-45-90 and 30-60-90 triangle relationships where needed.)

\medskip
\begin{center}
\renewcommand{\arraystretch}{2.4}
\begin{tabular}{|c|c|c|c|c|}
\hline
$\theta$ (radians) & $\theta$ (degrees) & $\sin\theta$ & $\cos\theta$ & $\tan\theta$ \\
\hline
$0$        & $0°$    & \blank{2cm} & \blank{2cm} & \blank{2cm} \\
$\pi/6$    & $30°$   & \blank{2cm} & \blank{2cm} & \blank{2cm} \\
$\pi/4$    & $45°$   & \blank{2cm} & \blank{2cm} & \blank{2cm} \\
$\pi/3$    & $60°$   & \blank{2cm} & \blank{2cm} & \blank{2cm} \\
$\pi/2$    & $90°$   & \blank{2cm} & \blank{2cm} & \blank{2cm} \\
\hline
\end{tabular}
\end{center}

\vspace{0.14in}

\textbf{2.} A point $P = (-0.28,\; 0.96)$ is on the unit circle.
State $\sin\theta$, $\cos\theta$, and $\tan\theta$ to two decimal places.

\vspace{0.50in}

\textbf{3.} Explain in a complete sentence why $\tan(\pi/2)$ is undefined.

\vspace{0.50in}

\textbf{4.} \textit{(AP-style multiple choice)} Which of the following could be the cosine of an
angle whose terminal ray is in Quadrant III?

\medskip
\begin{tabular}{lllll}
  (A) $\dfrac{3}{5}$ \qquad &
  (B) $-\dfrac{3}{5}$ \qquad &
  (C) $\dfrac{4}{5}$ \qquad &
  (D) $1$ &
\end{tabular}

\vspace{0.20in}

\textbf{5.} Given that $\cos\theta = -\dfrac{4}{5}$ and $\sin\theta > 0$, find $\tan\theta$.
Show all steps.

\vspace{0.65in}

\end{notesbox}

\begin{extensionbox}
\textbf{Extension: Deriving the Pythagorean Identity for Tangent}

We know $\sin^2\theta + \cos^2\theta = 1$.

\begin{enumerate}[leftmargin=1.8em, itemsep=4pt]
  \item Divide both sides of $\sin^2\theta + \cos^2\theta = 1$ by $\cos^2\theta$.
        Write out each term after dividing:

        \[
          \frac{\sin^2\theta}{\cos^2\theta} + \frac{\cos^2\theta}{\cos^2\theta}
          = \frac{1}{\cos^2\theta}
        \]

  \item Rewrite $\dfrac{\sin^2\theta}{\cos^2\theta}$ using the definition of tangent. \vspace{0.25in}

  \item State the resulting identity. \vspace{0.25in}
\end{enumerate}

\textit{This identity, $\tan^2\theta + 1 = \sec^2\theta$ (where $\sec\theta = 1/\cos\theta$),
will appear frequently in future units.}
\end{extensionbox}

\vspace{0.10in}

\begin{notesbox}{Preview: Lesson 3.3}
In \textbf{Lesson 3.3} we derive the exact coordinates of $P$ for
$\theta = \pi/6,\, \pi/4,\, \pi/3$ using 45-45-90 and 30-60-90 right-triangle ratios.
Your completed table in Problem 1 will be fully confirmed with geometric proofs.
\end{notesbox}

\end{document}
